%!TEX root = ../thesis.tex
%*******************************************************************************
%*********************************** First Chapter *****************************
%*******************************************************************************

\chapter{Introduction}

% **************************** Define Graphics Path **************************
\ifpdf
    \graphicspath{{Chapter1/Figs/Raster/}{Chapter1/Figs/PDF/}{Chapter1/Figs/}}
\else
    \graphicspath{{Chapter1/Figs/Vector/}{Chapter1/Figs/}}
\fi

\section{Background and Context}
Burnout among clinicians has become a well-documented challenge in modern medical practice. Nurses and doctors alike report high levels of exhaustion, impaired attentiveness, and a loss of meaning in their work, symptoms that have been linked to increased use of electronic health records (EHRs), administrative overload, and shift pressures \cite{nejm2016burnout, physicianburnout2018, timeallocation2016, burnoutcovid2021}. Psychiatric nurses, the focus of this thesis, combine clinical responsibilities with the emotionally demanding task of supporting individuals living with mental health conditions. Their day-to-day work requires continuous patient observation, medication management, coordination with multidisciplinary teams, and empathetic engagement with vulnerable patients \cite{understandingnursing2019}. These competing demands frequently leave little time for documentation or rapid access to the comprehensive patient information needed for safe care.

\section{Problem Statement and Research Gap}
Despite the promise of digital systems, most psychiatric clinics continue to rely on fragmented or handwritten patient records. Nurses often depend on memory, the single physical copy of the patient record, or time-consuming manual searches, diverting attention away from patient interaction. Existing conversational assistants primarily target patient-facing use cases or focus on medical specialties other than psychiatry, leaving a gap in tools that support psychiatric nurses in Greek clinical settings. At the same time, large language models (LLMs) and conversational AI have matured to the point where they can summarise and synthesise clinical information, yet their use raises questions about accuracy, privacy, and hallucinations \cite{llmsmodernmedicine2023, opportunitiesrisks2024, raglewisetal2021}. There is therefore a need to investigate whether an AI-powered assistant can reliably retrieve nurse-specific information while respecting the constraints of mental health care.

\section{Research Objectives and Questions}
The overarching objective of this research is to design, implement, and evaluate a conversational assistant that helps psychiatric nurses retrieve patient information on demand. The study addresses the following research questions:\label{chap:introduction-rqs}
\begin{itemize}
    \item RQ1) Is it feasible to develop a conversational assistant that helps nurses extract information about patients' profiles on demand?
    \item RQ2) If feasible, how effective is the assistant in terms of faithfulness, answer relevance and correctness, context precision, and context recall?
    \item RQ3) Does the assistant improve the daily experience and satisfaction of nurses when interacting with patient information?
\end{itemize}

\section{Significance of the Work}
Automating access to patient context through an AI assistant promises to reduce cognitive load, mitigate burnout, and allow nurses to focus on therapeutic interactions rather than administrative retrieval. Voice- and text-enabled conversational interfaces can support nurses who need hands-free access during shifts, provide structured summaries for rapid decision-making, and serve as a foundation for future integration with electronic systems. By grounding responses in curated patient data and highlighting limitations, the prototype contributes to the growing field of agentic AI in healthcare while foregrounding responsible deployment in high-stakes environments \cite{clinicalentityIE2025}.

\section{Structure of the Thesis}
This thesis is organised as follows. Chapter~\ref{chap:lit-review} surveys the literature on clinician burnout, healthcare conversational agents, and AI techniques relevant to psychiatric practice. Chapter~\ref{chap:methodology} details the methodology, including data sources, system design, tooling, and evaluation procedures. Chapter~\ref{chap:results-discussion} presents the experimental results and discusses their implications in relation to the research questions and prior work. Finally, Chapter~\ref{chap:conclusions} summarises the key findings, outlines the contributions, and proposes avenues for future research.

